\question[10]
Explain how the Coriolis effect and Ekman Transport together cause Antarctic upwelling. State specifically the direction (N, E, S, W) of the wind at this latitude and the direction of upwelled surface water flow.  The location of the Antarctic upwelling is about about 60$^\circ$S latitude.  You will find a global map on the last page of the exam for assistance.  You may draw on the map to indicate wind direction and water flow.

\vspace*{\stretch{4}}

\begin{solution}
Coriolis: 3 points
Ekman: 3 points
Wind: 2 points
\end{solution}

\question[2]
How many atmospheres of pressure does a squid living 560 meters deep in the Pacific Ocean experience?

\AnswerBox{\baselineskip}{%
	57 atmospheres.
}


%% Five point version with extra credit.
\question[5]
List the three deepest ocean trenches and the ocean basin in which each is located.  Then circle the deepest trench. Extra Credit: List the depth, within 100 feet / 30 meters, of each trench you list. 0.5 point per trench. Be sure to indicate whether you are using feet or meters.

\vspace{\baselineskip}

	\begin{tabular}{ccc}
		\large Trench Name		&	\large Ocean Basin	&	\large Depth \small(extra credit) \\[5ex]
		\makebox[5cm]{\hrulefill} & \makebox[5cm]{\hrulefill} & \makebox[4cm]{\hrulefill} \\[5ex]
		\makebox[5cm]{\hrulefill} & \makebox[5cm]{\hrulefill} & \makebox[4cm]{\hrulefill}\\[5ex]
		\makebox[5cm]{\hrulefill} & \makebox[5cm]{\hrulefill} & \makebox[4cm]{\hrulefill} \\
	\end{tabular}

%% Three point version, no extra credit.
\question[3]
List the three deepest ocean trenches and the ocean in which each trench is located.  1/2 point per blank. 

\vspace{0.5\baselineskip}

\begin{tabular}{cc}
	Trench Name		&	Ocean Basin \\[4ex]
	\makebox[5cm]{\ifprintanswers\textbf{Marianas/Challenger Deep}\else \hrulefill \fi} &
	\makebox[5cm]{\ifprintanswers\textbf{Pacific}\else \hrulefill \fi} \\[4ex]
	
	\makebox[5cm]{\ifprintanswers\textbf{Puerto Rico}\else\hrulefill \fi} &
	\makebox[5cm]{\ifprintanswers\textbf{Atlantic}\else \hrulefill \fi} \\[4ex]
	
	\makebox[5cm]{\ifprintanswers\textbf{Sunda/Diamantina}\else \hrulefill \fi} &
	\makebox[5cm]{\ifprintanswers\textbf{Indian}\else \hrulefill \fi} \\
\end{tabular}

%% Four point version, no extra credit. SHALLOWEST
\question[4]
List the three \emph{shallowest} ocean trenches and the ocean basin in which each is located. Circle the \emph{deepest} of the three. 0.5 point per blank plus 1 point for the circle. 

\vspace{\baselineskip}

\begin{tabular}{cc}
	\large Trench Name		&	\large Ocean Basin \\[5ex]
	\makebox[5cm]{\hrulefill} & \makebox[5cm]{\hrulefill} \\[5ex]
	\makebox[5cm]{\hrulefill} & \makebox[5cm]{\hrulefill} \\[5ex]
	\makebox[5cm]{\hrulefill} & \makebox[5cm]{\hrulefill} \\
\end{tabular}


\question[5]
Briefly describe one adaptation that allows a bottom-dwelling organism to live along a passive margin coastline.

\vspace*{\stretch{1}}

\question[5]
Briefly describe one adaptation that allows a bottom-dwelling organism to live along an active margin coastline.

\question[10]
Describe the difference between active and passive margins. Describe how this affects the biology of benthic (bottom-dwelling) organisms living on the different margin types. Be clear in your example. 



\vspace*{\stretch{1}}

\question[15]
Describe how the mechanism of plate tectonics causes the formation of new sea floor, the loss of older sea floor, and the movement of the plates (including the pushing and pulling forces). You must also include the types of crust involved, and their relative ages and density for a clear and thorough explanation. You may include a drawing. Use proper terminology. Do not tell me about earthquakes and volcanos. 

\question[15]
Explain slab pull and ridge push, and describe how the two together cause the movement of tectonic plates. Include the types of crust involved and relevant density information to support your answer. Use proper terminology. You may include an illustration. Do not tell me about earthquakes and volcanos. 


 	\begin{AnswerPage}{4\basespace}
	Ridge push: 
	\vspace*{0.5\baselineskip}
	
	1. Magma rises through the space at the edges of two oceanic crust plates, pushing up the plate edges. [2 points]
	\vspace*{0.5\baselineskip}
	
	2. The magma cools on the edge of the plates. [2 points]
	\vspace*{0.5\baselineskip}
	
	3. The weight of the plate edges pushes the plate away from the ridge.  [2 points]
	\vspace*{1\baselineskip}
	
	Slab pull:
	\vspace*{0.5\baselineskip}
	
	1. Denser oceanic crust meets and sinks below the less dense continental crust. [2 points]
	\vspace*{0.5\baselineskip}
	
	2. As oceanic crust sinks, it becomes even denser. [2 points]
	\vspace*{0.5\baselineskip}
	
	3. This causes it to sink faster, pulling along the rest of the plate. [2 points]
	\vspace*{1\baselineskip}
	
	[2 points allotted for clarity, proper terminology, grammar and spelling.]
	
	\end{AnswerPage}

\question
	\begin{parts}
	\part[6]
	On screen is a whipnose seadevil (genus \textit{Gigantactis}). State whether the whipnose seadevil more likely lives in the mesopelagic or bathypelagic zone, and explain one adaptation, other than bioluminescence and \emph{visible in the image}, the fish exhibits to justify your answer. The adaptation \emph{must be characteristic of the depth zone} and not common to all deep-sea fishes. Be sure to explain the specific relationship between the adaptation and the zone you choose.

		\begin{AnswerPage}{2\basespace}
			Bathypelagic [2 points], plus one of:
			\vspace*{0.5\baselineskip}

			\quad Small eyes [2 pts] don't expend energy for lightless zone but still see biolum. [2 points]
			\vspace*{0.5\baselineskip}

			\quad Blobby body {2 pts};  pressure provides support. No energy wasted on skeleton. [2 points]
			\vspace*{0.5\baselineskip}

			\quad Consider color. [2 points plus 2 points for explanation.]
		\end{AnswerPage}

	\part[5]
	Identify a different adaptation, other than bioluminescence and \emph{visible in the image}, that is specific to life in the deep-sea (this adaptation can be common to deep sea organisms regardless of depth zone). Explain how that adaptation is related to the environment or survival in the deep-sea. %Do not use bioluminescence.

		\begin{AnswerPage}{2\basespace}
			Apparently large stomach to hold large prey. [2 points] \\
			Nutrients are scarce. [3 points]
			\vspace*{0.5\baselineskip}

			Apparently large mouth to capture large prey. [2 points] \\
			Nutrients are scarce. [3 points]
			\vspace*{0.5\baselineskip}

			Consider color. [4 points]
			%Biolum lure to attract prey, which are scarce.
		\end{AnswerPage}

	\end{parts}


\question[6]
Briefly explain why the lecithotrophic larval strategy is more common than the planktotrophic strategy for marine organisms living at high latitudes. 

	\AnswerBox{\basespace}{%
	Planktotrophic larvae eat phytoplankton. The phytoplankton season is so short at high latitudes that a slight mistiming of larvae may miss the food supply. Lecithotrophic larvae do not feed so are not dependent on a food source so they have increased chance of survival.
}

\question[6]
Briefly explain how water temperature and salinity interact to cause thermohaline circulation in the North Atlantic Ocean. You do not need to explain the entire ocean conveyor.

\AnswerBox{\basespace}{%
	Water gets extremely cold in the North Atlantic. Cold water has increased density. Strong winds cause evaporation of the water. Salt does not evaporate, so the cold, salty water becomes even denser. It gets dense enough to sink below the deeper water.
}

\question
\begin{parts}
	\part[5]
	On screen is the Indo-Pacific Snaggletooth (\textit{Astronesthes indopacificus}). Tell whether the snaggletooth more likely lives in the mesopelagic or bathypelagic zone. Explain one adaptation, other than bioluminescence and \emph{visible in the image}, the fish exhibits to justify your answer. The adaptation \emph{must be characteristic of the depth zone} and not common to all deep-sea fishes. Be sure to explain the specific relationship between the adaptation and the zone you choose.
	
	\smallskip
	\emph{Write only one adaptation. If you write more than one, I will grade only the first one I read.}
	
	\begin{AnswerPage}{2\basespace}
		Mesopelagic [2 points], plus one of:
		\vspace*{0.5\baselineskip}
		
		\quad Large eyes [2 pts] to gather available light. [2 points]
		\vspace*{0.5\baselineskip}
		
		\quad Muscular body {2 pts};  for active swimming, migration to surface. [2 points]
		\vspace*{0.5\baselineskip}
		
		\quad Consider color. [2 points plus 2 points for explanation.]
	\end{AnswerPage}
	
	\part[5]
	Identify a different adaptation, other than bioluminescence and \emph{visible in the image}, that is specific to life in the deep-sea (this adaptation can be common to deep sea organisms regardless of depth zone). Explain how that adaptation is related to the environment or survival in the deep-sea. %Do not use bioluminescence.
	
	\smallskip
	\emph{Write only one adaptation. If you write more than one, I will grade only the first one I read.}
	
	\begin{AnswerPage}{2\basespace}
		Large teeth to capture andhold large prey. [2 points] \\
		Nutrients are scarce. [3 points]
		\vspace*{0.5\baselineskip}
		
		Apparently large mouth to capture large prey. [2 points] \\
		Nutrients are scarce. [3 points]
		\vspace*{0.5\baselineskip}
		
		Consider color. [4 points]
		%Biolum lure to attract prey, which are scarce.
	\end{AnswerPage}
	
\end{parts}

